% Section 4: CONTROL THEORY AND ALGORITHMS

\section{Control Theory and Algorithms}
\label{sec:control-theory}

This section covers the theoretical foundations and practical implementation of all control algorithms used in the VentCon2 system, including PID control, valve mapping, anti-windup protection, hysteresis compensation, and automatic tuning.

\subsection{PID Controller Theory}
\label{sec:pid-theory}

This section provides a thorough explanation of the PID (Proportional-Integral-Derivative) controller as used in the VentCon2 system. It covers the mathematical foundations, the two standard representations (parallel and standard/ISA form), the practical meaning of each tuning parameter, and the specific discrete-time implementation used by the firmware.

\subsubsection{The PID Control Law}

A PID controller continuously computes an \emph{error signal} $e(t)$ as the difference between a desired setpoint $r(t)$ and the measured process variable $y(t)$:

\begin{equation}
    e(t) = r(t) - y(t)
\end{equation}

In the VentCon2 system, $r(t)$ is the target pressure (in bar) set by the user, and $y(t)$ is the pressure measured by the sensor. The controller output $u(t)$ drives the proportional valve via PWM.

\subsubsection{Parallel (Independent) Form}
\label{sec:parallel-form}

The \textbf{parallel form} (also called the \emph{independent gains} or \emph{ideal} form) expresses the controller output as a direct sum of three independent terms:

\begin{equation}
    \boxed{u(t) \;=\; K_p \, e(t) \;+\; K_i \!\int_0^t e(\tau)\,d\tau \;+\; K_d \,\frac{de(t)}{dt}}
    \label{eq:pid-parallel}
\end{equation}

Each gain ($K_p$, $K_i$, $K_d$) has \textbf{different physical units} and acts completely independently:

\begin{table}[H]
\centering
\caption{Parallel-form PID parameters and their units}
\label{tab:pid-parallel-params}
\begin{tabular}{llll}
\toprule
\textbf{Parameter} & \textbf{Symbol} & \textbf{Units (VentCon2)} & \textbf{Role} \\
\midrule
Proportional gain & $K_p$ & PWM\,/\,bar & Instantaneous response to error \\
Integral gain     & $K_i$ & PWM\,/\,(bar\,$\cdot$\,s) & Accumulated error elimination \\
Derivative gain   & $K_d$ & PWM\,$\cdot$\,s\,/\,bar & Rate-of-change damping \\
\bottomrule
\end{tabular}
\end{table}

\noindent
\textbf{This is the form used by VentCon2.} The user-facing parameters $K_p$, $K_i$, $K_d$ entered via the web interface or serial commands are parallel-form gains. They are passed directly to the PID library without any conversion. The web interface additionally displays the equivalent standard-form time constants $T_i = K_p/K_i$ and $T_d = K_d/K_p$ for reference.

\paragraph{Advantages of the parallel form:}
\begin{itemize}
    \item Each term can be adjusted without affecting the others.
    \item Intuitive for manual tuning: change $K_i$ alone to fix steady-state offset.
    \item The auto-tuner (relay method with Ziegler--Nichols rules) directly outputs parallel gains.
\end{itemize}

\paragraph{Disadvantage:}
Because the three gains have different units and are independent, it can be harder to judge
the \emph{relative strength} of the I and D actions compared to the P action.

\subsubsection{Standard (ISA / Dependent) Form}

The \textbf{standard form} (ISA form, also called the \emph{dependent} or \emph{series} form in some references) factors out the proportional gain $K_p$ so that the integral and derivative actions are described by \emph{time constants}:

\begin{equation}
    \boxed{u(t) \;=\; K_p \!\left[\, e(t) \;+\; \frac{1}{T_i}\!\int_0^t e(\tau)\,d\tau \;+\; T_d \,\frac{de(t)}{dt} \,\right]}
    \label{eq:pid-standard}
\end{equation}

where:
\begin{itemize}
    \item $K_p$ is the proportional gain (same meaning as in the parallel form),
    \item $T_i$ is the \textbf{integral time} (also called \emph{reset time}), in seconds,
    \item $T_d$ is the \textbf{derivative time} (also called \emph{rate time}), in seconds.
\end{itemize}

\paragraph{Conversion between the two forms.}
The relationship is straightforward:

\begin{equation}
    K_i = \frac{K_p}{T_i}, \qquad K_d = K_p \cdot T_d
    \label{eq:form-conversion}
\end{equation}

\noindent
or equivalently:

\begin{equation}
    T_i = \frac{K_p}{K_i}, \qquad T_d = \frac{K_d}{K_p}
    \label{eq:form-conversion-inv}
\end{equation}

\begin{table}[H]
\centering
\caption{Comparison of standard vs.\ parallel form}
\label{tab:form-comparison}
\begin{tabular}{lcc}
\toprule
\textbf{Property} & \textbf{Parallel form} & \textbf{Standard (ISA) form} \\
\midrule
Controller equation     & $K_p\,e + K_i\!\int e\,dt + K_d\,\dot{e}$ & $K_p\bigl[e + \tfrac{1}{T_i}\!\int e\,dt + T_d\,\dot{e}\bigr]$ \\
Integral parameter      & $K_i$ \,[\,PWM/(bar$\cdot$s)\,] & $T_i$ \,[\,s\,] \\
Derivative parameter    & $K_d$ \,[\,PWM$\cdot$s/bar\,] & $T_d$ \,[\,s\,] \\
Gains independent?      & Yes & No ($T_i$, $T_d$ scale with $K_p$) \\
Changing $K_p$ affects   & P term only & \textbf{All three terms} \\
Used in VentCon2?       & \textbf{Yes} & No (but easy to convert) \\
\bottomrule
\end{tabular}
\end{table}

\subsubsection{Practical Meaning of the Time Constants}
\label{sec:time-constants-practical}

The time constants $T_i$ and $T_d$ have concrete physical interpretations that are useful for understanding controller behaviour, even though VentCon2 uses the parallel form internally.

\paragraph{Integral Time $T_i$ (Reset Time)}

$T_i$ answers the question: \emph{``If the error stayed constant, how long would it take for the integral action alone to produce the same output as the proportional action?''}

\begin{equation}
    \text{After time } T_i \text{ with constant error } e_0: \quad
    \underbrace{\frac{K_p}{T_i}\!\int_0^{T_i} e_0 \, d\tau}_{\text{integral contribution}} \;=\;
    \underbrace{K_p \, e_0}_{\text{proportional contribution}}
\end{equation}

\noindent
Practical implications:
\begin{itemize}
    \item \textbf{Small $T_i$} (e.g.\ 0.5\,s): Aggressive integral action. The controller quickly ramps up its output to eliminate even small offsets. \emph{Risk:} overshoot and oscillation from integral windup.
    \item \textbf{Large $T_i$} (e.g.\ 10\,s): Gentle integral action. The accumulated correction builds slowly. Steady-state error is eventually eliminated, but response is sluggish.
    \item \textbf{$T_i = \infty$} (i.e.\ $K_i = 0$): No integral action at all. A permanent steady-state offset will remain if the P term alone cannot drive the error to zero. This is the default factory setting in VentCon2.
\end{itemize}

\noindent
\textbf{Rule of thumb:} Start with $T_i$ roughly equal to the system's dominant time constant (the time the pressure takes to reach about 63\,\% of its final value after a step change in valve opening).

\paragraph{Derivative Time $T_d$ (Rate Time)}

$T_d$ answers the question: \emph{``How far into the future does the derivative action effectively `predict' the error based on its current rate of change?''}

\begin{equation}
    \text{D-term output} = K_p \, T_d \, \frac{de}{dt}
    \quad\Longleftrightarrow\quad
    \text{``P-term output that \emph{would} occur after } T_d \text{ seconds if the current trend continued''}
\end{equation}

\noindent
Practical implications:
\begin{itemize}
    \item \textbf{Small $T_d$} (e.g.\ 0.01\,s): Minimal damping effect. The controller barely reacts to the rate of change of the error signal. Useful when sensor noise is significant.
    \item \textbf{Moderate $T_d$} (e.g.\ 0.1--1\,s): Effective damping. The controller ``sees'' the pressure approaching setpoint and starts backing off early, reducing overshoot. This is the primary purpose of the D term.
    \item \textbf{Large $T_d$} ($>$\,1\,s): Very aggressive rate response. Highly sensitive to noise; even small pressure fluctuations cause large output swings. Rarely useful without additional filtering.
    \item \textbf{$T_d = 0$} (i.e.\ $K_d = 0$): No derivative action. The controller only reacts to the current error and its history, not its trend. This is the default factory setting in VentCon2.
\end{itemize}

\noindent
\textbf{Rule of thumb:} $T_d$ is typically set to $\frac{1}{4}\,T_i$ to $\frac{1}{8}\,T_i$ (for many process types). If the pressure signal is noisy, reduce $T_d$ or increase sensor filtering first.

\paragraph{Numerical Example (VentCon2 context)}

Suppose a user has tuned the system with the following parallel-form gains:
\[
K_p = 500, \quad K_i = 200, \quad K_d = 50
\]

The equivalent standard-form time constants are:
\[
T_i = \frac{K_p}{K_i} = \frac{500}{200} = 2.5\;\text{s}, \qquad
T_d = \frac{K_d}{K_p} = \frac{50}{500} = 0.1\;\text{s}
\]

\noindent
Interpretation:
\begin{itemize}
    \item With a constant pressure error, the integral action will match the proportional action's contribution after \textbf{2.5\,seconds}.
    \item The derivative action provides damping equivalent to predicting the error \textbf{0.1\,seconds} into the future.
    \item The ratio $T_d / T_i = 0.04$ indicates moderate derivative action relative to integral --- appropriate for a system where some overshoot damping is desired without excessive noise sensitivity.
\end{itemize}


\subsubsection{Discrete-Time Implementation in VentCon2}
\label{sec:pid-discrete}

The PID library used by VentCon2 (\texttt{Arduino-PID-Library}, imax9000 fork) implements a discrete-time approximation of Equation~\ref{eq:pid-parallel}. The library pre-scales the user-supplied gains by the sample time at configuration time, so each \texttt{Compute()} call only needs simple multiply-accumulate operations, which is efficient on the ESP32.

\paragraph{Gain pre-scaling.}
When \texttt{SetTunings($K_p$, $K_i$, $K_d$)} is called with sample time $T_s$ (in seconds):
\begin{align}
    k_p &= K_p               & &\text{(unchanged)} \label{eq:kp-internal} \\
    k_i &= K_i \cdot T_s     & &\text{(scaled for summation per sample)} \label{eq:ki-internal} \\
    k_d &= K_d \,/\, T_s     & &\text{(scaled for finite difference per sample)} \label{eq:kd-internal}
\end{align}

These internal values $k_p$, $k_i$, $k_d$ are the ones actually used in the compute loop.

\paragraph{Compute algorithm.}
Each sample period, the library executes (in Proportional-on-Error mode, which is the VentCon2 default):

\begin{align}
    e[n]          &= r[n] - y[n]                         & &\text{(error)} \\
    \Delta y[n]   &= y[n] - y[n{-}1]                     & &\text{(input change, not error change)} \\
    \mathit{outputSum} &\mathrel{+}= k_i \cdot e[n]      & &\text{(integral accumulator)} \\
    u[n]          &= k_p \cdot e[n] + \mathit{outputSum} - k_d \cdot \Delta y[n]
                                                          & &\text{(total output)}
\end{align}

\noindent
Key implementation details:

\begin{enumerate}
    \item \textbf{Derivative on measurement:} The derivative term uses $\Delta y$ (change in the measured pressure) rather than $\Delta e$ (change in error). This avoids the ``derivative kick'' that would occur when the setpoint changes abruptly --- a step change in setpoint would cause a momentary spike in $de/dt$, but $dy/dt$ remains smooth.
    
    \item \textbf{Integral via running sum:} The integral is computed as a running accumulation (\texttt{outputSum}) rather than storing and summing the entire error history. This is both memory-efficient and numerically stable.
    
    \item \textbf{Output clamping:} Both \texttt{outputSum} and the final output are clamped to $[0,\, \mathit{pwmFullScaleRaw}]$ (0 to 16383 for 14-bit PWM resolution), providing basic integral anti-windup at the library level.
    
    \item \textbf{Sample time enforcement:} \texttt{Compute()} checks elapsed time and only recalculates when $\geq T_s$ has passed. VentCon2 configures $T_s = 10\;\text{ms}$ (i.e.\ 100\,Hz PID update rate) via \texttt{SetSampleTime()}.
    
    \item \textbf{Control loop vs. PID sample rate:} The FreeRTOS control task runs at 1000\,Hz by default, but the PID library internally skips computation if less than $T_s$ has elapsed. This means the sensor is read at 1\,kHz (providing better filtering resolution), but the PID output updates at 100\,Hz.
\end{enumerate}

\subsubsection{Initialization and Startup Behaviour}
\label{sec:pid-init}

On power-up, the PID controller is initialized in the following sequence:

\begin{enumerate}
    \item Settings are loaded from flash (\texttt{/settings.json} on LittleFS), providing the persisted $K_p$, $K_i$, $K_d$, setpoint, sample time, and PWM resolution.
    \item \texttt{SetOutputLimits(0, 16383)} configures the valid output range.
    \item \texttt{SetSampleTime(10)} sets the PID sample period to 10\,ms and internally rescales $k_i$, $k_d$.
    \item \texttt{SetTunings($K_p$, $K_i$, $K_d$)} applies the loaded gains with correct scaling.
    \item One sensor reading is taken so that \texttt{pressureInput} contains the actual ambient pressure.
    \item \texttt{SetMode(Automatic)} transitions the PID from Manual to Automatic, triggering the internal \texttt{Initialize()} function which performs a \emph{bumpless transfer}: \texttt{outputSum} is set to the current output value (0), and \texttt{lastInput} is set to the current pressure, preventing a derivative spike on the first cycle.
    \item Only after full PID configuration are the FreeRTOS control tasks started.
\end{enumerate}

\noindent
\textbf{Factory default gains are $K_p = K_i = K_d = 0$.} This means the controller produces zero output on a fresh device --- the valve remains closed until the user configures non-zero gains. This is a deliberate safety choice: shipping with non-zero gains on an uncharacterized pressure system could cause dangerous overshoot or oscillation.

\subsubsection{Summary: Choosing Between Forms}

\begin{itemize}
    \item When \textbf{reading tuning literature} or controller datasheets, check whether they use the standard form ($K_p$, $T_i$, $T_d$) or the parallel form ($K_p$, $K_i$, $K_d$) --- the parameters look similar but have different units and cannot be swapped directly.
    \item When entering values into VentCon2's web interface: supply \textbf{parallel-form gains} ($K_p$, $K_i$, $K_d$). If a tuning guide gives standard-form constants, convert using Equation~\ref{eq:form-conversion}: $K_i = K_p / T_i$ and $K_d = K_p \cdot T_d$.
    \item The auto-tuner built into VentCon2 outputs parallel-form gains directly --- no conversion is needed.
\end{itemize}

\subsection{PWM Valve Control and Output Mapping}
\label{sec:valve-mapping}

\subsubsection{Physical Valve Characteristics}
The proportional valve used in the VentCon2 system does not respond linearly across the full 0--100\% PWM duty cycle range. Two physical constraints define the valve's effective operating region:

\begin{itemize}
    \item \textbf{Below $\sim$20\% duty cycle:} The valve coil does not generate sufficient electromagnetic force to overcome the return spring and internal friction. PWM energy in this range is dissipated as heat without producing valve movement.
    \item \textbf{Above $\sim$90\% duty cycle:} The valve plunger is fully retracted and the flow path is completely open. Additional current produces no further displacement and only generates excess heat in the coil.
\end{itemize}

These limits are defined in \texttt{Constants.h}:

\begin{lstlisting}[caption=Valve Operating Range Constants]
namespace ValveConfig {
    constexpr float PID_MIN_OUTPUT_PERCENT = 1.0f;  // Dead zone threshold
    constexpr float VALVE_MIN_DUTY = 20.0f;         // Minimum effective duty cycle
    constexpr float VALVE_MAX_DUTY = 90.0f;         // Maximum useful duty cycle
}
\end{lstlisting}

\subsubsection{Output Mapping Function}
To ensure the PID controller has full resolution over the valve's useful range, the system maps the PID's logical 0--100\% output to the valve's physical 20--90\% operating range. This mapping is performed by the \texttt{mapPIDoutputToPwmValve()} function in \texttt{ControlSystem.cpp}:

\begin{lstlisting}[caption=PID Output to Valve PWM Mapping Function]
uint32_t ControlSystem::mapPIDoutputToPwmValve(double pidOutput, int maxPwmFullScaleRaw) {
    // Convert PID output to percentage (0-100%)
    float pidPercent = (pidOutput / maxPwmFullScaleRaw) * 100.0;
    
    // If below minimum threshold, keep valve closed
    if (pidPercent < ValveConfig::PID_MIN_OUTPUT_PERCENT) {
        return 0;
    }
    
    // Map PID's 0-100% to valve's effective range
    float mappedPercent = ValveConfig::VALVE_MIN_DUTY + (pidPercent / 100.0) * 
                         (ValveConfig::VALVE_MAX_DUTY - ValveConfig::VALVE_MIN_DUTY);
    
    // Constrain to valid range
    mappedPercent = constrain(mappedPercent, 0.0, 100.0);
    
    // Convert percentage back to absolute PWM value
    return (uint32_t)((mappedPercent / 100.0) * maxPwmFullScaleRaw);
}
\end{lstlisting}

The mapping follows this equation:

\begin{equation}
D_{valve} = D_{min} + \frac{PID_{\%}}{100} \cdot (D_{max} - D_{min})
\end{equation}

where $D_{valve}$ is the actual valve duty cycle, $D_{min} = 20\%$, $D_{max} = 90\%$, and $PID_{\%}$ is the PID output as a percentage of full scale. For 14-bit PWM resolution ($PWM_{fullscale} = 16383$), the final raw PWM value is:

\begin{equation}
PWM_{raw} = \frac{D_{valve}}{100} \cdot PWM_{fullscale}
\end{equation}

\subsubsection{Mapping Table}
The following table illustrates representative mapping values for 14-bit resolution:

\begin{center}
\begin{tabular}{r r r r}
\hline
\textbf{PID Output (\%)} & \textbf{PID Raw Value} & \textbf{Valve Duty (\%)} & \textbf{PWM Raw Value} \\
\hline
$<$ 1\% & $<$ 164 & 0\% & 0 (closed) \\
1\% & 164 & 20.7\% & 3\,391 \\
25\% & 4\,096 & 37.5\% & 6\,144 \\
50\% & 8\,192 & 55.0\% & 9\,011 \\
75\% & 12\,287 & 72.5\% & 11\,878 \\
100\% & 16\,383 & 90\% & 14\,745 \\
\hline
\end{tabular}
\end{center}

\subsubsection{Dead Zone Threshold Design}
A critical design decision is the choice of the dead zone threshold at 1\% (\texttt{PID\_MIN\_OUTPUT\_PERCENT}) rather than at the 20\% valve minimum. This threshold determines when the valve transitions from fully closed (0\% duty) to minimum opening (20\% duty). The 1\% value is chosen deliberately for control-theoretic reasons:

\begin{itemize}
    \item \textbf{Minimized dead time:} When the PID controller determines that the valve should open (output $>$ 0), it only needs to cross the 1\% threshold before the valve responds at 20\% duty. If the threshold were set at 20\%, the PID would need to ramp its output through the entire 0--20\% range---a region where the valve is physically unresponsive. During this ramp, the PID receives no pressure feedback, causing the integral term to accumulate unchecked.
    
    \item \textbf{Reduced integral windup:} With a 20\% threshold, the PID integral would grow continuously while the output traverses the 0--20\% dead zone. Even with anti-windup protection resetting the integrator at the boundaries, there is a timing window where the PID crosses the threshold with significant accumulated momentum, leading to overshoot when the valve finally opens.
    
    \item \textbf{Immediate feedback:} The 1\% threshold ensures that the transition from ``valve closed'' to ``valve open'' occurs almost as soon as the PID requests any output. The valve immediately jumps to 20\% duty---the minimum at which it physically responds---and the PID receives pressure feedback within one control cycle. This tight feedback loop enables precise regulation from the moment the valve opens.
    
    \item \textbf{Intentional discontinuity:} The mapping creates a deliberate step from 0\% to 20\% duty at the 1\% threshold. This discontinuity mirrors the physical reality of the valve: there is no useful intermediate state between ``energized below opening force'' and ``energized at minimum opening force.'' The mapping makes the control signal match the actuator's actual behavior.
\end{itemize}

\subsubsection{Interaction with Anti-Windup}
The valve mapping and anti-windup mechanisms work together to maintain control stability. The anti-windup logic in the control loop checks whether the PID output percentage falls below \texttt{VALVE\_MIN\_DUTY} (20\%) or above \texttt{VALVE\_MAX\_DUTY} (90\%). When the output is in these saturation regions and moving further into saturation, the PID integrator is reset to prevent windup. The mapping function then translates the constrained PID output into the appropriate PWM value. This two-stage approach---first limiting integral accumulation, then mapping to the physical range---ensures robust control across the entire operating envelope.

\subsection{Anti-Windup Protection}
\label{sec:anti-windup}

\subsubsection{The Integral Windup Problem}
Integral windup (also called integrator saturation) is a critical issue in PID control systems where the integral term accumulates error continuously even when the system cannot respond further. This occurs when:

\begin{itemize}
    \item The actuator (valve) reaches its physical limits (fully open or closed)
    \item The PID output saturates at its maximum or minimum values
    \item The error persists but the system cannot reduce it further
\end{itemize}

In the VentCon2 system, this manifests as follows: When the proportional valve reaches its maximum duty cycle and pressure continues to rise above setpoint, the integral term continues to accumulate the error. When the setpoint is finally reduced or the pressure drops, the accumulated integral error causes the output to remain high longer than necessary, resulting in:

\begin{itemize}
    \item \textbf{Overshoot:} The system overshoots the new setpoint by a significant margin
    \item \textbf{Oscillation:} The system exhibits prolonged oscillations around the target
    \item \textbf{Sluggish Response:} Slow recovery from disturbances
    \item \textbf{Poor Stability:} Difficulty achieving steady-state control
\end{itemize}

Mathematically, the integral term continues to accumulate:

\begin{equation}
I = \int_0^t e(\tau) \, d\tau
\end{equation}

Even when the output is saturated and cannot change:

\begin{equation}
u(t) = K_p e(t) + K_i \int_0^t e(\tau) \, d\tau + K_d \frac{de(t)}{dt}
\end{equation}

If $u(t) > u_{max}$, the system is saturated, but the integral continues to grow, leading to a large $I$ term that persists even after the error direction reverses.

\subsubsection{Anti-Windup Implementation in VentCon2}
The VentCon2 system implements anti-windup protection through conditional integral reset when output saturation is detected. This is controlled exclusively via the \textbf{AW ON} and \textbf{AW OFF} serial commands; no web interface toggle is currently available for this feature.

The anti-windup mechanism works as follows:

\begin{lstlisting}[caption=Anti-Windup Logic from ControlSystem.cpp]
// Anti-windup for deadband and saturation
if (settings->antiWindup) {
    float pidPercent = (*pwmPIDoutput / *pwmFullScaleRaw) * 100.0;
    
    // Check if output is saturating
    if ((pidPercent < ValveConfig::VALVE_MIN_DUTY && *pwmPIDoutput < previousOutput) ||
        (pidPercent > ValveConfig::VALVE_MAX_DUTY && *pwmPIDoutput > previousOutput)) {
        // Reset the PID to prevent integral accumulation
        pid->SetMode(PID::Manual);
        pid->SetMode(PID::Automatic);
    }
}
\end{lstlisting}

The algorithm detects saturation in two scenarios:

\begin{itemize}
    \item \textbf{Lower Saturation:} Output is below minimum valve duty ($pidPercent < 20\%$) and attempting to increase
    \item \textbf{Upper Saturation:} Output is above maximum valve duty ($pidPercent > 90\%$) and attempting to increase further
\end{itemize}

When saturation is detected, the system performs a PID reset by toggling between manual and automatic modes, which clears the accumulated integral error and re-initializes the controller state. This prevents the integral term from growing unbounded when the system cannot respond.

\subsubsection{Configuration and Tuning}
Anti-windup can be controlled via:

\begin{itemize}
    \item \textbf{Serial Command:} \texttt{AW ON} / \texttt{AW OFF}
    \item \textbf{Default State:} Disabled by default for systems that don't require it
    \item \textbf{Persistent:} Setting is automatically saved to flash memory
\end{itemize}

\textbf{Note:} Anti-windup is not currently exposed in the web interface and must be configured via serial commands.

For pressure control applications where the valve has significant deadband or limited operating range, enabling anti-windup is recommended to improve response stability and reduce oscillation.

\subsection{Hysteresis Compensation}
\label{sec:hysteresis}

\subsubsection{The Hysteresis Problem}
Hysteresis is a lag or delay in the system's response to changes in the control signal. In proportional valve systems, several sources of hysteresis exist:

\begin{itemize}
    \item \textbf{Magnetic Hysteresis:} The valve coil exhibits nonlinear behavior with hysteresis in the B-H curve
    \item \textbf{Mechanical Friction:} Static friction in the valve plunger requires a minimum force to initiate movement
    \item \textbf{Pressure Drop Effects:} Pressure differential across the valve affects its opening threshold
    \item \textbf{Thermal Lag:} Temperature changes in the valve coil alter its magnetic properties
    \item \textbf{Flow Rate Effects:} The relationship between PWM and actual flow is nonlinear
\end{itemize}

Hysteresis causes several control problems:

\begin{itemize}
    \item \textbf{Deadband:} A region where small control signal changes produce no output change
    \item \textbf{Hunting:} Continuous oscillation around the setpoint as the system alternates between on and off
    \item \textbf{Limit Cycles:} Persistent oscillation with amplitude and frequency determined by the deadband size
    \item \textbf{Poor Steady-State:} Inability to achieve precise pressure regulation
\end{itemize}

The effect can be modeled as:

\begin{equation}
y(t) = \begin{cases}
y_{min} & \text{if } u(t) < u_{on} \\
y(t-\Delta t) & \text{if } u_{off} < u(t) < u_{on} \\
y_{max} & \text{if } u(t) > u_{off}
\end{cases}
\end{equation}

Where $u_{on} > u_{off}$ represents the deadband width, and the output depends on where it came from (history-dependent).

\subsubsection{Hysteresis Compensation Strategy}
The VentCon2 system implements hysteresis compensation by adding a dead zone or bias to the PID output that adjusts based on the current system state. This is enabled via the \textbf{HYST ON} / \textbf{HYST OFF} commands and configured with \textbf{HYSTAMT <value>}.

The compensation mechanism works by introducing an offset to the setpoint or error signal that creates a predictable margin:

\begin{itemize}
    \item When pressure is below setpoint and rising: Apply a positive offset to assist opening
    \item When pressure is above setpoint and falling: Apply a negative offset to assist closing
    \item When approaching setpoint: Reduce the offset smoothly to allow fine control
\end{itemize}

This is implemented in the control loop as:

\begin{lstlisting}[caption=Hysteresis Compensation Mechanism]
// Hysteresis compensation to reduce hunting around setpoint
if (settings->hysteresis) {
    float pressureDifference = *pressureInput - lastPressure;
    float compensationAmount = (settings->hystAmount / 100.0) * settings->setpoint;
    
    // Adjust setpoint based on pressure trend
    if (pressureDifference > 0) {
        // Pressure rising: reduce effective setpoint to encourage closing
        adjustedSetpoint = settings->setpoint - compensationAmount;
    } else if (pressureDifference < 0) {
        // Pressure falling: increase effective setpoint to encourage opening
        adjustedSetpoint = settings->setpoint + compensationAmount;
    } else {
        // Stable: use exact setpoint
        adjustedSetpoint = settings->setpoint;
    }
    
    // Use adjusted setpoint for PID error calculation
    error = adjustedSetpoint - *pressureInput;
}
\end{lstlisting}

The compensation creates a virtual deadband around the setpoint:

\begin{equation}
\text{Deadband} = \frac{\text{HYSTAMT}}{100} \cdot \text{Setpoint}
\end{equation}

For example, with a 3 bar setpoint and HYSTAMT = 5\%:
$$\text{Deadband} = 0.05 \times 3 = 0.15 \text{ bar}$$

This creates an effective hysteresis band of ±0.15 bar around the setpoint, reducing the system's tendency to hunt.

\subsubsection{Configuration and Tuning}
Hysteresis compensation can be controlled via:

\begin{itemize}
    \item \textbf{Enable/Disable:} \texttt{HYST ON} / \texttt{HYST OFF} serial commands
    \item \textbf{Compensation Amount:} \texttt{HYSTAMT <value>} (0-20\%, default 5\%)
    \item \textbf{Typical Range:} 2-10\% for most applications
    \item \textbf{Persistent Storage:} Automatically saved to flash
\end{itemize}

\textbf{Note:} Like anti-windup, hysteresis compensation is not currently exposed in the web interface and must be configured via serial commands.

\subsubsection{Tuning Guidelines}

The following guidelines help optimize hysteresis compensation for specific applications:

\begin{itemize}
    \item \textbf{High Sensitivity Requirement:} Use lower HYSTAMT (2-3\%) for applications requiring tight control
    \item \textbf{Noisy Sensor:} Use higher HYSTAMT (8-10\%) if sensor readings are noisy
    \item \textbf{Sluggish Response:} Reduce HYSTAMT if the system responds too slowly to setpoint changes
    \item \textbf{Oscillation:} Increase HYSTAMT if the system oscillates around the setpoint
    \item \textbf{Interaction with Filter:} Hysteresis works best when combined with appropriate filter strength (FLT command)
\end{itemize}

\subsubsection{Interaction with Other Control Parameters}

The effectiveness of hysteresis compensation depends on the configuration of other parameters:

\begin{itemize}
    \item \textbf{Low-Pass Filter:} The filter strength (FLT 0.0-1.0) smooths sensor noise and works synergistically with hysteresis
    \item \textbf{PID Gains:} Proportional gain (Kp) affects sensitivity; higher Kp requires lower HYSTAMT
    \item \textbf{Valve Deadband:} Physical valve deadband (VALVE\_MIN\_DUTY and VALVE\_MAX\_DUTY) adds to effective hysteresis
    \item \textbf{Control Frequency:} Faster control loops (higher CONTROL FREQ) can reduce hunting effect
\end{itemize}

A typical tuning sequence would be:

\begin{enumerate}
    \item Enable both anti-windup and hysteresis compensation via serial commands (\texttt{AW ON}, \texttt{HYST ON})
    \item Set HYSTAMT to 5\% as a starting point (\texttt{HYSTAMT 5})
    \item Adjust FLT (filter strength) to 0.2-0.5 to smooth sensor readings (via web interface or \texttt{FLT 0.3})
    \item Perform auto-tuning (\texttt{TUNE START}) to get optimal PID gains
    \item Fine-tune HYSTAMT based on observed oscillation:
    \begin{itemize}
        \item If still hunting: increase HYSTAMT by 1-2\%
        \item If sluggish response: decrease HYSTAMT by 1-2\%
    \end{itemize}
\end{enumerate}

\subsection{Relay-Based Auto-Tuning}
\label{sec:auto-tuning}

The VentCon2 system includes an automatic PID tuning capability that uses relay-based oscillation to identify the critical frequency and amplitude of the controlled system, then applies established tuning rules to calculate optimal PID parameters.

\subsubsection{Tuning Algorithm}
The auto-tuner uses relay oscillation to identify the critical frequency and amplitude of the system, then applies tuning rules to calculate optimal PID parameters.

\begin{equation}
K_c = \frac{4M}{\pi A}
\end{equation}

\begin{equation}
T_c = 2 \cdot T_{osc}
\end{equation}

Where:
\begin{itemize}
    \item $K_c$ = Critical gain
    \item $M$ = Relay amplitude  
    \item $A$ = Process oscillation amplitude
    \item $T_c$ = Critical period
    \item $T_{osc}$ = Measured oscillation period
\end{itemize}

\subsubsection{Tuning Rules}
\begin{table}[H]
\centering
\begin{tabular}{@{}llll@{}}
\toprule
\textbf{Rule} & \textbf{Kp} & \textbf{Ki} & \textbf{Kd} \\
\midrule
Ziegler-Nichols Classic & $0.6 K_c$ & $\frac{2K_p}{T_c}$ & $\frac{K_p T_c}{8}$ \\
Ziegler-Nichols Aggressive & $0.33 K_c$ & $\frac{2K_p}{T_c}$ & $\frac{K_p T_c}{3}$ \\
Tyreus-Luyben & $0.454 K_c$ & $\frac{K_p}{2.2 T_c}$ & $\frac{K_p T_c}{6.3}$ \\
Pessen Integral & $0.7 K_c$ & $\frac{2.5 K_p}{T_c}$ & $\frac{K_p T_c}{6.25}$ \\
\bottomrule
\end{tabular}
\caption{Auto-Tuning Rules Implemented}
\end{table}

During auto-tuning, the system operates in manual PWM mode and switches the valve on/off within the configured PWM range to induce oscillation. After sufficient oscillation cycles are observed, the algorithm calculates the critical gain and period, then applies the selected tuning rule to compute PID gains. The user can accept or reject the calculated parameters via serial commands.
